\section{Discussion and future work}
\subsection{Scope of the current conclusions}
The formulation and search apply to general simple undirected graphs.
This is distinct from experimental scope: the $L(3,2)$ tree screen and
multi-family pilot are reported separately because their graph domains
and budgets differ. In the pilot, Gurobi solves more instances to optimality
than CaDiCaL, while one instance is solved to optimality only by CaDiCaL.
Encoding correctness does not depend on which backend is faster. Trees
with 1600 vertices provide an initial size study; short runtimes are not
an error and do not establish that trees are easy in general. In this
sample, 48 trees attain the degree bound and two show that it is not always tight.

Only cycles use root fixing in the symmetry experiments. Reflections on
product families require explicit automorphisms, not merely equal vertex
degrees. Reducing symmetric solutions does not imply proportional reductions
in clauses or runtime. Families without additional symmetry constraints
must be described as repeated runs of the same model. These results are
not an evaluation of incremental SAT: learned clauses are not reused across spans.

\subsection{Limitations}
A valid labeling establishes an upper bound. Optimality additionally requires
a matching lower bound, a proof excluding smaller spans, or an exact formula
within its validity domain. Agreement between two backends is a cross-check,
not a general proof. A value between theoretical lower and upper bounds is
consistent with them even if it attains neither. Finite agreement with a
formula is validation evidence, not a proof over an entire family.

Greedy labeling currently supplies incumbents and MIP starts; it has not
been timed as a separate heuristic baseline. Heuristics and metaheuristics
may find good solutions quickly but do not automatically prove optimality.
A solution attaining a proven lower bound can nevertheless be certified
optimal. We neither assume that every heuristic is faster than exact methods
nor infer optimality from solution quality alone. A separate heuristic
benchmark requires common label-domain conventions, validation, and timing definitions.
A GA adapted from radio labeling \cite{sarhan2026genetic} should use a common
budget and multiple algorithmic seeds, together with a random-order multi-start
baseline to isolate the benefit of evolution. Its best solution remains an
upper bound unless supported by additional optimality evidence. No GA data
are currently available to compare quality or speed with SAT--ILP.

The new screen does not include a direct-versus-order ablation, specialized
tree algorithms, or MILP formulations other than assignment. Results for
one formulation/backend cannot establish superiority over all ILP methods.
Graph-generation seeds measure input variation; repeated runs on one graph
measure timing variation. These are not interchangeable. Short runtimes
also reflect startup and licensing overhead; each comparison must use a
consistent timing definition.

Known formulas and the reference proofs in the appendix validate the
pipeline and are not claimed as new mathematical contributions. The lower
bound $m+4$ for $C_n\circ P_m$ does not establish equality without an upper-bound
construction. A survey note or AI-generated proposal does not replace an
original publication when establishing novelty or complexity for a graph class.

\subsection{Next priorities}
The repeated confirmation experiment has now been completed. The immediate
priority is to review the formula/bound references within their applicability
domains and interpret instances for which optimality remains unresolved.
Results favoring Gurobi must be retained; the dataset should not be reselected
to increase the fraction of SAT wins. The SAT--ILP comparison should be
consolidated before proceeding to the proposed GA baseline.
A subsequent SAT improvement should be evaluated by ablation, for example
solver reuse across spans or more efficient constraint generation. Independently
checkable UNSAT certificates would strengthen optimality conclusions where
mathematical bounds are insufficient. Multiple techniques should not be
introduced simultaneously and their combined gains attributed to a single technique.

Further evaluation should distinguish model-construction cost from search
cost and retain instances not solved within budget. This helps identify
where an improvement acts, rather than inferring its mechanism solely
from total runtime or clause counts at one span.
